\documentclass[11pt,a4paper]{article}
\usepackage{fontspec,amsmath,amssymb,amsthm,microtype}
\setmainfont{lmroman10-regular.otf}[BoldFont=lmroman10-bold.otf,ItalicFont=lmroman10-italic.otf,BoldItalicFont=lmroman10-bolditalic.otf]
\setmonofont{lmmono10-regular.otf}[Scale=MatchLowercase]
\usepackage[margin=25mm]{geometry}
\usepackage[unicode,hidelinks]{hyperref}
\setlength{\emergencystretch}{3em}
\setlength{\parskip}{0.4em}
\setlength{\parindent}{0pt}
\allowdisplaybreaks[2]
\raggedbottom
\title{\texorpdfstring{Subgaussian Gibbs score and its sharp second moment}{Subgaussian Gibbs score and its sharp second moment}}
\author{LangevinLean mathematical exposition}
\date{}
\begin{document}\maketitle
\noindent Source module: \texttt{Unadjusted\allowbreak{}Langevin/\allowbreak{}Gibbs\allowbreak{}Score.\allowbreak{}lean}. The accompanying HTML pairs this mathematical exposition with the exact Lean statements and records the scope of the formalization.


Differentiability, an \(L\)-Lipschitz gradient, and finite Gibbs normalization suffice for \(\int\|\nabla U\|^2d\pi\le Ld\). Strong convexity is used only by the final wrapper to derive normalization.


\section{\texorpdfstring{gibbs\allowbreak{}Weight mul exp score le}{gibbsWeight mul exp score le}}


The descent lemma gives \(U(x+v)\le U(x)+\langle\nabla U(x),v\rangle+(L/2)\|v\|^2\). Negate and exponentiate, then rearrange to \(e^{-U(x)}e^{-\langle\nabla U(x),v\rangle}\le e^{L\|v\|^2/2}e^{-U(x+v)}\). Lean first establishes the real inequality and then embeds it in the extended nonnegative numbers.



{\raggedright\noindent Lean declarations: \texttt{gibbs\allowbreak{}Weight\_\allowbreak{}mul\_\allowbreak{}exp\_\allowbreak{}score\_\allowbreak{}le}.\par}

\section{\texorpdfstring{gibbs score lmgf le}{gibbs score lmgf le}}


Integrate the pointwise exponential-score bound. Translation invariance gives \(\int e^{-U(x+v)}dx=Z_U\); normalization cancels this factor, leaving \(\mathbb E_{\pi_U}e^{-\langle\nabla U,v\rangle}\le e^{L\|v\|^2/2}\). Finiteness of \(Z_U\) is explicit because extended inverse cancellation needs it.



{\raggedright\noindent Lean declarations: \texttt{gibbs\_\allowbreak{}score\_\allowbreak{}lmgf\_\allowbreak{}le}.\par}

\section{\texorpdfstring{gibbs score subgaussian}{gibbs score subgaussian}}


Apply the preceding bound with the vector \(-tv\), for every real \(t\). The exponent becomes \(t\langle\nabla U,v\rangle\) and the right side is \(e^{t^2L\|v\|^2/2}\). This is exactly \texttt{Has\allowbreak{}Subgaussian\allowbreak{}MGF} with variance proxy \(L\|v\|^2\); the proof also establishes exponential integrability required by that structure.



{\raggedright\noindent Lean declarations: \texttt{gibbs\_\allowbreak{}score\_\allowbreak{}subgaussian}.\par}

\section{\texorpdfstring{sq le exp add exp neg sub two}{sq le exp add exp neg sub two}}


The series \(\cosh a=\sum_{n\ge0}a^{2n}/(2n)!\) has nonnegative terms, so retaining its first two terms gives \(\cosh a\ge1+a^2/2\). Substitute \(2\cosh a=e^a+e^{-a}\) and rearrange to \(a^2\le e^a+e^{-a}-2\). Lean uses \texttt{Real.\allowbreak{}has\allowbreak{}Sum\_\allowbreak{}cosh} with \texttt{sum\_\allowbreak{}le\_\allowbreak{}has\allowbreak{}Sum}, explicitly proving nonnegativity of every discarded even-power term.



{\raggedright\noindent Lean declarations: \texttt{sq\_\allowbreak{}le\_\allowbreak{}exp\_\allowbreak{}add\_\allowbreak{}exp\_\allowbreak{}neg\_\allowbreak{}sub\_\allowbreak{}two}.\par}

\section{\texorpdfstring{integral sq le of subgaussian}{integral sq le of subgaussian}}


Integrate \(t^2X^2\le e^{tX}+e^{-tX}-2\), giving \(t^2\mathbb EX^2\le2(e^{ct^2/2}-1)\). Subgaussianity provides the exponential integrability and hence the square integrability needed here. Divide by \(t^2\) for \(t\ne0\), and let \(t\to0\) to obtain \(\mathbb EX^2\le c\). Lean computes the scalar limit before applying closedness of order.



{\raggedright\noindent Lean declarations: \texttt{integral\_\allowbreak{}sq\_\allowbreak{}le\_\allowbreak{}of\_\allowbreak{}subgaussian}.\par}

\section{\texorpdfstring{integral norm sq gradient gibbs le}{integral norm sq gradient gibbs le}}


Choose an orthonormal basis \((e_i)_{i=1}^d\) and use Parseval: \(\|\nabla U(x)\|^2=\sum_i\langle\nabla U(x),e_i\rangle^2\). Each coordinate is subgaussian with proxy \(L\), hence has second moment at most \(L\). Finite-sum integrability permits integration term by term, yielding \(\mathbb E_{\pi_U}\|\nabla U\|^2\le Ld\).



{\raggedright\noindent Lean declarations: \texttt{integral\_\allowbreak{}norm\_\allowbreak{}sq\_\allowbreak{}gradient\_\allowbreak{}gibbs\_\allowbreak{}le}.\par}

\section{\texorpdfstring{mem\allowbreak{}Lp gradient gibbs\allowbreak{}Measure}{memLp gradient gibbsMeasure}}


The coordinate subgaussian bounds prove integrability of every squared gradient coordinate. Summing over the finite orthonormal basis proves integrability of \(\|\nabla U\|^2\). Lipschitz continuity provides strong measurability, and the \texttt{Mem\allowbreak{}Lp} square criterion packages those facts as \(\nabla U\in L^2(\pi_U)\).



{\raggedright\noindent Lean declarations: \texttt{mem\allowbreak{}Lp\_\allowbreak{}gradient\_\allowbreak{}gibbs\allowbreak{}Measure}.\par}

\section{\texorpdfstring{integral norm sq gradient gibbs le of strong\allowbreak{}Convex}{integral norm sq gradient gibbs le of strongConvex}}


Strong convexity with \(m>0\) supplies the finite-partition premise of the score theorem. Substitute that proof into the first-order smoothness result to obtain \(\mathbb E_{\pi_U}\|\nabla U\|^2\le Ld\). The strong convexity constant is used for normalization here; it does not appear in the sharp score bound.



{\raggedright\noindent Lean declarations: \texttt{integral\_\allowbreak{}norm\_\allowbreak{}sq\_\allowbreak{}gradient\_\allowbreak{}gibbs\_\allowbreak{}le\_\allowbreak{}of\_\allowbreak{}strong\allowbreak{}Convex}.\par}

\end{document}
